\documentclass[a4paper,12pt]{article}

\usepackage[T1]{fontenc}
\usepackage[italian]{babel}
\usepackage{graphicx}
\usepackage{geometry}
\usepackage{tabularx}
\usepackage[table]{xcolor}
\usepackage[most]{tcolorbox}
\usepackage{fancyhdr}
\usepackage[hidelinks]{hyperref}

\newcommand{\groupmail}{alittlebyte.swe@gmail.com}
\newcommand{\grouplogo}{../../../.github/site-src/assets/logo_alittlebyte.png}
\newcommand{\groupsymbol}{../../../.github/site-src/assets/solo_logo.png}

\geometry{
    left=2.5cm,
    right=2.5cm,
    top=2.5cm,
    bottom=2.5cm,
    headheight=331pt,
    includeheadfoot
}

\pagestyle{fancy}
\fancyhf{}

\renewcommand{\headrulewidth}{0pt}
\fancyhead{}

\renewcommand{\footrulewidth}{0.4pt}
\fancyfoot[L]{\includegraphics[height=0.6cm]{\groupsymbol}\hspace{0.45em}aLittleByte}
    \fancyfoot[C]{\thepage}
\fancyfoot[R]{Verbale 2026-07-29}

\fancypagestyle{firstpage}{
    \fancyhf{}
    \renewcommand{\headrulewidth}{0pt}
    \renewcommand{\footrulewidth}{0.4pt}
    \fancyhead[C]{%
        \makebox[\textwidth][c]{%
            \begin{tabular}{c}
                \includegraphics[height=10cm]{\grouplogo}\\[1ex]
                \small\texttt{\groupmail}
            \end{tabular}%
        }
    }
    \fancyfoot[L]{\includegraphics[height=0.6cm]{\groupsymbol}\hspace{0.45em}aLittleByte}
        \fancyfoot[C]{\thepage}
    \fancyfoot[R]{Verbale 2026-07-29}
}

\begin{document}

\thispagestyle{firstpage}

\vspace*{0.5cm}

\begin{center}
    {\LARGE \textbf{Verbale Interno - 2026-07-29}}
\end{center}

\vspace{1.5cm}

\begin{center}
    \renewcommand{\arraystretch}{1.3}
    \begin{tcolorbox}[
        width=0.92\textwidth,
        colback=gray!6,
        colframe=gray!45,
        boxrule=0.5pt,
        arc=2mm,
        boxsep=0pt,
        left=8pt, right=8pt, top=8pt, bottom=8pt
    ]
        \begin{tabularx}{\linewidth}{>{\bfseries}p{3.5cm} X}
            Gruppo       & aLittleByte (gruppo 19) \\
            Responsabile & Alex Oloni \\
            Verificatore & Angelica Gastaldello\\
            Destinatari  & Prof. Tullio Vardanega, Prof. Riccardo Cardin \\
            Data         & 29 Luglio 2026\\
        \end{tabularx}
    \end{tcolorbox}
\end{center}

\newpage

\newgeometry{left=2.5cm,right=2.5cm,top=2.5cm,bottom=2.5cm,includefoot}
\tableofcontents
\newpage
\section{Informazioni Generali}
\section*{Specifiche Documento}
\hrule
\vspace{1cm}

\begin{tabular}{>{\bfseries}p{5cm} | p{10cm}}
Luogo Riunione & Discord \\[6pt]

Orario (Inizio - Fine) & 10:30 -- 11:15 \\[6pt]

Responsabile &
    A. Oloni\\[6pt]

Amministratore & 
    A. Maule\\[6pt]

Verificatore &
    A. Gastaldello\\[6pt]

Partecipanti &
  A. Oloni\\ &
  L. Soligo\\&
  E. Bellet\\&
  A. Gastaldello\\&
  A. Maule\\&
  V. Y. Kaneda Fernandes\\&
  D. Belluz\\

\end{tabular}


\section{Ordine del Giorno}
    \begin{itemize}
    \item Avanzamento del progetto e analisi delle funzionalità mancanti
    \item Gestione repository: allineamento branch develop
    \item Assegnazione Use Case 21, 22, 39.12 e 39.15
    \item Gestione riuso della configurazione e Use Case 32
    \end{itemize}


\section{Attività svolte}
\subsection{Avanzamento del progetto}
In primis, il team ha discusso dello stato attuale di avanzamento del progetto, facendo il punto della situazione sulle attività completate e su quelle ancora mancanti per raggiungere i prossimi obiettivi prefissati.

\subsection{Gestione repository e branch develop}
Successivamente si è discusso dell'allineamento del codice. Si è deciso che Leonardo Soligo si occuperà di raccogliere i lavori svolti ed effettuare il merge di tutto il codice aggiornato all'interno del branch \texttt{develop}.

\subsection{Sviluppo Use Case e assegnazione task}
Il team ha poi esaminato i prossimi casi d'uso (Use Case) da implementare, suddividendo il lavoro come segue:
\begin{itemize}
    \item \textbf{UC 21 e 22:} L'implementazione di questi due casi d'uso è stata affidata a Eleonora Bellet.
    \item \textbf{UC 39.12 e 39.15:} Diego Belluz prenderà in carico questo blocco. Nello specifico, si dovrà occupare delle funzionalità relative all'inserimento della mail e alla modifica dei campi destinatario, data e ora.
    \item \textbf{Configurazione e UC 32:} Angela Maule e Angelica Gastaldello lavoreranno congiuntamente all'implementazione del caso d'uso 32 e alla logica di riuso delle configurazioni.
\end{itemize}


\newpage
\section{To-Do List}
\begin{table}[h]
    \centering
    \begin{tabular}{|p{4.5cm}|p{5cm}|l|}
    \hline
        \rowcolor{lightgray}\textit{\textbf{Persona Incaricata}}  & \textbf{Task Assegnata} &\textbf{Link Issue} \\
    \hline
        \textit {Leonardo Soligo} & Inserimento modifiche e merge nel branch develop & \href{https://github.com/aLittleByte-19/Documentazione/issues/157}{\#157}\\
    \hline
        \textit {Eleonora Bellet} & Implementazione UC 21 e UC 22 & \href{https://github.com/aLittleByte-19/Documentazione/issues/158}{\#158}\\
    \hline
        \textit{Diego Belluz} & Implementazione UC 39.12, 39.15 (mail, destinatario, data, ora) & \href{https://github.com/aLittleByte-19/Documentazione/issues/159}{\#159}\\
    \hline
        \textit{Angela Maule e Angelica Gastaldello} & Riuso della configurazione e sviluppo UC 32 & \href{https://github.com/aLittleByte-19/Documentazione/issues/160}{\#160}\\
    \hline
        \textit{Diego Belluz} & Stesura Verbale Interno 29/07/2026 & \href{https://github.com/aLittleByte-19/Documentazione/issues/161}{\#161}\\
    \hline
    \end{tabular}
    \label{tab:placeholder}
\end{table}

\end{document}